\chapter{Plan de estudio por bloques}\label{app:study}
La guía fija 150 horas totales para 6 ECTS. La distribución siguiente es una propuesta editorial que incluye asistencia, estudio y evaluación; no modifica horarios ni la distribución oficial de actividades. Cada bloque combina lectura activa, derivación sin consultar la fórmula final, resolución y corrección de errores.
\begin{center}\begin{tabular}{llr}\toprule Bloque&Objetivo de salida&Horas orientativas\\\midrule
I&Balances, unidades y prestaciones&20\\
II&Tobera completa y selección de régimen&30\\
III&Geometría, presión y límites de estabilidad&20\\
IV&Mezcla, térmica, alimentación y termoquímica&35\\
V&Potencia, haces y consumo eléctrico&20\\
VI&Etapas, trayectoria y maniobra&25\\\midrule Total&&150\\\bottomrule\end{tabular}\end{center}
\section{Primera vuelta: construir modelos}
Para cada capítulo, escribir el volumen de control, las conservaciones y las hipótesis antes de memorizar resultados. Reproducir al menos un despeje sin mirar la solución. En toberas, dibujar rama y contrapresión; en sólidos, elegir la variable geométrica; en líquidos, cerrar las dos corrientes; en eléctrico, fijar frontera de potencia; en vehículo, escribir todas las masas presentes.
\section{Segunda vuelta: resolver y comprobar}
Resolver los ejercicios propuestos antes de leer el apéndice H. Utilizar después los problemas BASIC e INTERMEDIATE para consolidar conversiones y relaciones. Marcar por separado errores de unidades, signos, modelo, álgebra y aritmética. Repetir sólo la familia que produjo el error, con datos distintos, y comprobar que la corrección se sostiene por un balance independiente.
\section{Tercera vuelta: integración y examen}
Trabajar ADVANCED y EXAM\_LEVEL explicando el régimen y las limitaciones, además de los resultados. Realizar un simulacro en 150 minutos con su hoja de datos. Corregir con soluciones separadas y rúbrica. En otra sesión, repetir sin tiempo las partes falladas y justificar cada hipótesis. La distribución de cinco familias en los simulacros es editorial y no una predicción del examen oficial.
\section{Trabajo y defensa}
Reservar dentro del tiempo de estudio la preparación del trabajo y su defensa, que representan el 25\,\% según la guía. Deben realizarse también para la convocatoria correspondiente según las condiciones docentes. La prueba escrita representa el 75\,\%; la guía puede exigir condiciones por componente sin fijar aquí un mínimo numérico universal. Preparar una explicación oral de balances, fuentes, modelo y límites ayuda tanto al trabajo como a los problemas.
\section{Lista de dominio antes de la evaluación}
Ser capaz de derivar empuje y Tsiolkovski; calcular caudal y prestaciones sin imponer adaptación; comprobar primer bloqueo y choque; distinguir geometría y tiempo en regresión; cerrar masa y potencia de alimentación; convertir correctamente energía, masa molar e Isp; y explicar qué resultados son ideales. Una cifra sin hipótesis ni unidad no demuestra dominio aunque coincida con una solución.
